% Generated draft: requires mathematical and source review.
\documentclass[11pt]{article}
\usepackage[margin=25mm]{geometry}
\usepackage{fontspec}
\setmainfont{Latin Modern Roman}
\usepackage{amsmath,amssymb,mathtools,unicode-math}
\setmathfont{Latin Modern Math}
\usepackage{enumitem,needspace,xurl,hyperref}
\setlength{\parindent}{0pt}
\setlength{\parskip}{5pt}
\setlength{\emergencystretch}{3em}
\newcommand{\N}{\mathbb{N}}
\newcommand{\Z}{\mathbb{Z}}
\newcommand{\Q}{\mathbb{Q}}
\newcommand{\R}{\mathbb{R}}
\newcommand{\C}{\mathbb{C}}
\newcommand{\F}{\mathbb{F}}
\newcommand{\A}{\mathbb{A}}
\newcommand{\PP}{\mathbb{P}}
\newcommand{\E}{\mathbb{E}}
\newcommand{\eps}{\varepsilon}
\newcommand{\dd}{\,\mathrm{d}}
\begin{document}
\section*{UnsolvedMath ID: 30003433}
\textbf{Source problem code:} catalogue entry
\subsection{Definitions and mathematical statement}
The supplied problem asks for the best possible degree bounds in sum-of-squares representations of nonnegative polynomials and in associated rational-function representations. The coefficient field, domain of nonnegativity, permitted certificates, and degree convention are not specified. Consequently, a unique precise optimization problem cannot be recovered from the record.

One possible interpretation, conditional on global nonnegativity over the real numbers and rational sum-of-squares certificates, can be formulated as follows. For integers \(n\geq 1\) and \(d\geq 0\), let
\[
\mathcal P_{n,d}=\{f\in\R[x_1,\ldots,x_n]:\deg f\leq d,\ f(a)\geq0\text{ for every }a\in\R^n\}.
\]
Here degree means total degree, with the zero polynomial regarded as satisfying every nonnegative degree bound. A common-denominator rational sum-of-squares certificate for \(f\) consists of an integer \(m\geq1\) and polynomials \(q,p_1,\ldots,p_m\in\R[x_1,\ldots,x_n]\), with \(q\neq0\), such that
\[
q^2f=\sum_{i=1}^{m}p_i^2.
\]
Equivalently, this is the identity
\[
f=\sum_{i=1}^{m}\left(\frac{p_i}{q}\right)^2
\quad\text{in }\R(x_1,\ldots,x_n).
\]
The polynomial identity holds everywhere; the displayed quotients are pointwise defined only where \(q\neq0\).

Under this interpretation, define the optimal uniform bound by
\[
B(n,d)=\inf\left\{D\in\Z_{\geq0}:\begin{array}{l}
\text{for every }f\in\mathcal P_{n,d}\text{ there exist }m\geq1,\ q\neq0,\ p_1,\ldots,p_m\\
\text{in }\R[x_1,\ldots,x_n]\text{ such that }q^2f=\sum_{i=1}^{m}p_i^2,\\
\deg q\leq D\text{ and }\deg p_i\leq D\text{ for every }1\leq i\leq m
\end{array}\right\},
\]
with \(\inf\varnothing=+\infty\). The corresponding conditional problem is to determine \(B(n,d)\), or sharp bounds for its dependence on \(n\) and \(d\). This formulation leaves the number of squares unrestricted.

\emph{Unresolved scope.} The source does not establish that this particular interpretation is intended. An identity \(f=\sum_i p_i^2\) using polynomials alone is a different requirement and is not available for every globally nonnegative real polynomial. Rational certificates, polynomial identities after clearing denominators, and certificates on specified semialgebraic sets therefore require separate conventions. The source's claimed equivalence does not specify how degree bounds are transferred between representations. No defining inequalities for a semialgebraic set, coefficient restrictions, restrictions on denominators, or bound on the number of squares are supplied.
\subsection{Short English statement}
Determine sharp degree bounds for sum-of-squares certificates of nonnegative polynomials. The source leaves the field, positivity domain, certificate type, and meaning of degree unspecified. A conditional global real formulation above asks for the least uniform bound on numerator and common-denominator degrees in terms of the number of variables and the degree of the input polynomial; the intended formulation requires human review.
\subsection{Sources}
\begin{itemize}
\item Local UnsolvedMath registry: supplied complete record, permanent ID \(30003433\), source problem code \texttt{catalogue entry}, title ``Optimal Degree Bounds for Sum-of-Squares Representations,'' in the Oberwolfach Reports Open Problems collection. No individual proposer is identified. The supplied proposed year is 2017.
\item Supplied report citation: \emph{Real Algebraic Geometry With a View Toward Moment Problems and Optimization} (2018), \emph{Oberwolfach Reports}, pp.\ 771--862, DOI \texttt{10.4171/owr/2017/14}. \url{https://doi.org/10.4171/owr/2017/14}. The record attributes the problem to workshop open-problem discussions; no exact problem page or speaker is supplied.
\item Supplied literature reference: Lombardi--Perrucci--Roy, ``effective Hilbert seventeenth problem.'' \url{https://arxiv.org/abs/1404.2338}.
\item Supplied literature reference: ``Representations of positive polynomials on semialgebraic sets.'' \url{https://doi.org/10.1016/j.jmaa.2006.05.028}.
\item Literature assessment dated \texttt{2026-08-21}: the supplied curation labels the problem partially solved and reports constructive elementary-recursive degree bounds for Hilbert's seventeenth problem and the Positivstellensatz, together with lower bounds in special settings, while stating that best possible general bounds remain unresolved. This is supplied dated curation, not an independently verified literature assessment; its status and cited claims require independent verification.
\item The registry supplies the statement label \texttt{corrected\_verified} and says that mathematical content was preserved while notation, typography, and sentence structure were repaired for publication. These are supplied editorial assertions, not independent mathematical verification of this draft.
\end{itemize}
\end{document}
